\documentclass{article}
\usepackage{blindtext}
\usepackage[a4paper, total={7in, 9.2in}]{geometry}
\usepackage{titling}
\usepackage{tabularx}
\usepackage{hyperref}
\usepackage{array}
\usepackage{longtable}
\usepackage[table]{xcolor}
\usepackage{booktabs}
\usepackage{float}
\usepackage{fancyhdr}

\pagestyle{fancy}
\fancyhf{}
\renewcommand{\headrulewidth}{0.4pt}
\fancyhead[C]{\textsc{The National University of Computer and Emerging Sciences, Lahore}}
\fancyfoot[C]{\thepage}

\title{
{\Large\scshape The National University of Computer and Emerging Sciences, Lahore}\\[10pt]
{\large Course Outline}\\[4pt]
{\Huge\textbf{Compiler Construction CS-4031}}\\[8pt]
{\normalsize Semester Fall-2026}
}
\author{ }
\date{}

\pretitle{\begin{center}}
\posttitle{\end{center}}
\preauthor{}
\postauthor{}
\predate{}
\postdate{}

\begin{document}

\maketitle
\vspace{-0.5cm}

\begin{table}[H]
    \centering
    \def\arraystretch{1.5}
    \begin{tabularx}{\textwidth}{ l X l X }
        \textbf{Coordinator}: & Dr. Faisal Aslam & \textbf{Office hours}: & Monday to Thursday \\
        & & & 10:00 to 11:00\\
        \textbf{Credit Hours:} & 3 & \textbf{Prerequisites:} & Programming Fundamentals\\
        & & & Theory of Automata \\
        \multicolumn{4}{l}{\textbf{Instructors:} Mr. Aamir Raheem, Mr. Hamad ul Qudous, Dr. Faisal Aslam, Dr. Farooq Ahmed Rana} \\
    \end{tabularx}
    \label{tab:Intro}
\end{table}
\vspace{-1cm}

\section*{Course Description}
Compiler Construction introduces students to the theory and practice of translating high-level programs into executable code. The course follows a working compiler front-to-back: lexical analysis, syntax analysis (parsing), semantic analysis, intermediate-code generation, run-time environments, and code optimization. Emphasis is placed on the formal models that make each phase tractable -- regular expressions and finite automata for scanning, context-free grammars and both top-down and bottom-up parsing (including the industry-standard LALR(1) algorithm) for parsing, and data-flow analysis for optimization -- alongside hands-on construction of each phase using standard tools (Flex, Yacc/Bison) and a semester-long project. By the end of the course, students should be able to explain why each phase of a compiler is structured the way it is, analyze and choose appropriate parsing and optimization algorithms for a given language, and independently design and implement the lexical analyzer, parser, and intermediate-code generator for a small language.

\section*{Program Learning Outcomes (PLOs)}
This course covers the following PLOs:

\renewcommand{\arraystretch}{1.4}
\begin{longtable}{|l|p{3.3cm}|p{10.3cm}|}
\hline
\textbf{PLO\#}  & \textbf{PLO Name} & \textbf{PLO Description} \\ \hline
\endfirsthead
\hline
\textbf{PLO\#}  & \textbf{PLO Name} & \textbf{PLO Description} \\ \hline
\endhead
PLO 2 & Knowledge for Solving Computing Problems &
Apply knowledge of computing fundamentals, knowledge of a computing specialization, and mathematics, science, and domain knowledge appropriate for the computing specialization to the abstraction and conceptualization of computing models from defined problems and requirements \\ \hline
PLO 3 & Problem Analysis & Identify, formulate, research literature, and solve complex computing problems reaching substantiated conclusions using fundamental principles of mathematics, computing sciences, and relevant domain disciplines \\ \hline
PLO 4 & Design/Development of Solutions & Design and evaluate solutions for complex computing problems, and design and evaluate systems, components, or processes that meet specified needs with appropriate consideration for public health and safety, cultural, societal, and environmental considerations \\ \hline
PLO 5&  Modern Tool Usage & Focuses on creating, selecting, adapting, and applying modern tools, techniques, and resources to solve complex computing problems, understanding the limitations of these tools \\ \hline
PLO 6 & Individual and Team Work & Aims to help students function effectively both as individuals and as team members or leaders in diverse and multidisciplinary settings \\ \hline
\end{longtable}

\section*{Course Learning Outcomes (CLOs)}
CS-4031 is a core Computer Science course with Programming Fundamentals and Theory of Automata as prerequisites. The course learning outcomes are:

\renewcommand{\arraystretch}{1.4}
\begin{longtable}{|l|p{9cm}|p{2.7cm}|p{2.2cm}|}
\hline
\textbf{CLO \#}  & \textbf{CLO Description} & \textbf{BT Level} & \textbf{PLO \#}\\ \hline
\endfirsthead
\hline
\textbf{CLO \#}  & \textbf{CLO Description} & \textbf{BT Level} & \textbf{PLO \#}\\ \hline
\endhead
CLO 1 & Describe the basic phases of compilation and recognize the theories and techniques used in each phase  & Understand (BT Level 2) & PLO 2 \\ \hline
       CLO 2 & Analyze and recognize the significance of the several phases through which a typical program is compiled, including a state-of-the-art parsing algorithm (LALR(1)) and a code-optimization algorithm (common subexpression elimination via data-flow analysis) & Analyze (BT Level 4) & PLO 4, PLO 3 \\ \hline
CLO 3 & Design and implement a lexical analyzer, parser, and intermediate-code generator for a given language, culminating in a semester project & Create (BT Level 6) & PLO 6, PLO 4, PLO 5 \\ \hline
\end{longtable}

\section*{Primary Textbooks}
These are cited by section number throughout the lecture table below and should be treated as the core references for lecture preparation.
\begin{enumerate}
    \item[{[DB]}] \textit{Compilers: Principles, Techniques, and Tools} (2nd ed.) by Aho, Lam, Sethi, Ullman --- \textbf{the "Dragon Book"; primary reference for section numbers below}
    \item[{[ET]}] \textit{Engineering a Compiler} (2nd ed.) by Cooper and Torczon
    \item[{[AP]}] \textit{Modern Compiler Implementation in C/Java/ML} by Appel
\end{enumerate}

\section*{Supplementary / Further Reading}
Referenced occasionally for specific topics; useful for deeper treatment but not required for lecture preparation.
\begin{enumerate}
    \item[{[PP]}] \textit{Programming Language Pragmatics} by Scott
    \item[{[MU]}] \textit{Advanced Compiler Design and Implementation} by Muchnick
    \item[{[AN]}] \textit{The Definitive ANTLR 4 Reference} by Parr
    \item[{[GR]}] \textit{Compilers: A Survey} by Gries
\end{enumerate}

\section*{Online / Tool References}
\begin{enumerate}
    \item[{[LLVM-LR]}] LLVM Language Reference Manual \\
    \href{https://llvm.org/docs/LangRef.html}{\small\texttt{llvm.org/docs/LangRef.html}}
    \item[{[LLVM-KA]}] "My First Language Frontend with LLVM" (Kaleidoscope Tutorial) \\
    \href{https://llvm.org/docs/tutorial/}{\small\texttt{llvm.org/docs/tutorial}}
    \item[{[ANTLR]}] ANTLR4 project site \& documentation \\
    \href{https://www.antlr.org/}{\small\texttt{antlr.org}}
    \item[{[TS]}] Tree-sitter --- incremental, error-tolerant parsing library \\
    \href{https://tree-sitter.github.io/tree-sitter/}{\small\texttt{tree-sitter.github.io/tree-sitter}}
\end{enumerate}


\section*{Lecture-Wise Course Outline (30 Lectures)}
\renewcommand{\arraystretch}{1.35}
\setlength{\tabcolsep}{4pt}
\footnotesize
\begin{longtable}{|c|p{5.1cm}|p{1.1cm}|p{2.5cm}|p{3.3cm}|}
\hline
\textbf{L\#} & \textbf{Topic} & \textbf{CLO} & \textbf{Dragon Book [DB]} & \textbf{Other References} \\ \hline
\endfirsthead
\hline
\textbf{L\#} & \textbf{Topic} & \textbf{CLO} & \textbf{Dragon Book [DB]} & \textbf{Other References} \\ \hline
\endhead

\multicolumn{5}{|c|}{\cellcolor{gray!20}\textbf{Unit 1: Introduction and Motivation (2 Lectures)}} \\ \hline
1 & Compilers vs. interpreters vs. JIT; why study compilers; applications of compiler technology & CLO1 & \S1.1, \S1.4, \S1.5 & [ET] \S1.1; [AP] Ch.1 \\ \hline
2 & Structure of a compiler: phases overview; front-/middle-/back-end; preprocessor, assembler, linker, loader & CLO1 & \S1.2, \S2.1--\S2.8 (overview) & [PP] \S1.5, \S1.6 \\ \hline

\multicolumn{5}{|c|}{\cellcolor{gray!20}\textbf{Unit 2: Lexical Analysis (5 Lectures)}} \\ \hline
3 & Tokens, patterns, lexemes; lexical errors; role of the lexical analyzer & CLO1 & \S3.1, \S3.3 & [AP] Ch.2 \\ \hline
4 & Regular languages vs. non-regular languages; regular expressions and regular definitions, including advanced operators, with extra worked examples & CLO1 & \S3.3, \S3.4 & [PP] \S3.3 \\ \hline
5 & Finite automata; transition diagrams; RE $\to$ NFA (Thompson's Construction) & CLO1, CLO2 & \S3.6, \S3.7 (Alg.~3.23) & [AP] \S2.3--2.4 \\ \hline
6 & NFA $\to$ DFA (subset construction); DFA minimization via transition diagrams; NFA vs.\ DFA power/performance & CLO1, CLO2 & \S3.7, \S3.9 (Alg.~3.20, 3.36) & [ET] \S2.4--2.5 \\ \hline
7 & Building a lexical analyzer using a lexer-generator tool (hands-on) & CLO3 & \S3.5, \S3.8 & [AP] \S2.5 \\ \hline
\rowcolor{blue!8}
Opt.1 & \emph{\textbf{[Optional -- instructor's tool choice: Flex]}} Building the L7 lexer concretely with Flex. & CLO3 & (tool choice for \S3.5, \S3.8) & --- \\ \hline

\multicolumn{5}{|c|}{\cellcolor{gray!20}\textbf{Unit 3a: Context-Free Grammars (Lectures 8--10)}} \\ \hline
8 & Context-free grammars; derivations; parse trees; ambiguity & CLO1 & \S4.1--\S4.2, \S4.3.3 & [GR] Ch.4 \\ \hline
9 & Grammar engineering I: eliminating left recursion and left factoring & CLO1, CLO2 & \S4.3 & [ET] \S3.5 \\ \hline
10 & Grammar engineering II: removing ambiguity (e.g.\ dangling-else), disambiguation techniques, non-CFG constructs & CLO1, CLO2 & \S4.3.3, \S4.3 & [GR] Ch.4 \\ \hline

\multicolumn{5}{|c|}{\cellcolor{gray!20}\textbf{Unit 3b: Parsing Algorithms (Lectures 11--17)}} \\ \hline
11 & Top-down parsing overview; recursive-descent parsing & CLO1, CLO3 & \S4.4.1 & [AP] \S3.2 \\ \hline
\multicolumn{5}{|c|}{\cellcolor{yellow!25}\textit{--- Midterm 1 (after Lecture 10; covers L1--L10) ---}} \\ \hline
12 & FIRST and FOLLOW sets; LL(1) parsing-table construction & CLO1, CLO2 & \S4.4.2--\S4.4.4 & [ET] \S3.3 \\ \hline
13 & LL(1) predictive-parsing algorithm; panic-mode error recovery & CLO2, CLO3 & \S4.4.3, \S4.4.5--\S4.4.9 & [PP] \S3.4.5 \\ \hline
14 & Bottom-up parsing: shift-reduce parsing, handles, shift-reduce \& reduce-reduce conflicts & CLO1, CLO2 & \S4.5 & [AP] \S3.3 \\ \hline
15 & LR(0) items \& automaton; SLR parsing & CLO2 & \S4.6 & [ET] \S3.4.1--3.4.2 \\ \hline
16 & \textbf{LALR(1) parsing} --- the industry-standard, state-of-the-art table-driven algorithm (used by Yacc/Bison), with a brief comparative look at canonical LR(1) & CLO2 & \S4.7.1--\S4.7.5 (canonical LR(1) covered briefly; LALR(1) in full) & [ET] \S3.4.3--3.4.4; [AN] Ch.1 \\ \hline
17 & Using ambiguous grammars in practice; building a parser with a parser-generator tool (hands-on) & CLO3 & \S4.8, \S4.9 & [AN] Ch.4--5 \\ \hline
\rowcolor{blue!8}
Opt.2 & \emph{\textbf{[Optional -- instructor's tool choice: Bison/Yacc]}} Building the L17 parser concretely with Bison/Yacc. & CLO3 & (tool choice for \S4.8, \S4.9) & --- \\ \hline
\rowcolor{blue!8}
Opt.3 & \emph{\textbf{[Optional -- modern alternative]}} ANTLR4: generates a combined lexer + adaptive LL(*) parser from one grammar file, with strong IDE/multi-language support -- a direct modern successor to Lex/Yacc. & CLO3 & (not in DB; extends \S4.9) & [ANTLR] ANTLR4 site \\ \hline
\rowcolor{blue!8}
Opt.4 & \emph{\textbf{[Optional -- modern alternative]}} Tree-sitter: incremental, error-tolerant GLR-style parsing -- the current state of the art for editor/IDE-grade parsing (used by Neovim, GitHub, VS Code). & CLO3 & (not in DB; extends \S4.9) & [TS] Tree-sitter site \\ \hline

\multicolumn{5}{|c|}{\cellcolor{gray!20}\textbf{Unit 4a: Syntax-Directed Definitions (Lecture 18)}} \\ \hline
18 & Syntax-directed definitions; S-attributed and L-attributed grammars & CLO1 & \S5.1--\S5.2 & [ET] Ch.4 intro \\ \hline
19 & \textbf{Review 2:} parsing algorithms (top-down, bottom-up, LR family, LALR(1)) \& intro SDT -- consolidation and worked examples (no new material) & CLO1, CLO2 & (review of \S4.4--4.9, \S5.1--5.2) & --- \\ \hline

\multicolumn{5}{|c|}{\cellcolor{gray!20}\textbf{Unit 4b: Semantic Analysis (Lectures 20--21)}} \\ \hline
20 & Syntax trees; SDT schemes; introduction to type checking & CLO1, CLO2 & \S5.3--\S5.5, \S6.5 & [PP] \S6.1--6.2 \\ \hline
\multicolumn{5}{|c|}{\cellcolor{yellow!25}\textit{--- Midterm 2 (after Lecture 20; covers L11--L20) ---}} \\ \hline
21 & Symbol tables, scope resolution, static semantic checking, with emphasis on practical implementation & CLO1, CLO2 & \S2.7, \S6.3 & [AP] Ch.5 \\ \hline

\multicolumn{5}{|c|}{\cellcolor{gray!20}\textbf{Unit 5: Intermediate Code Generation (4 Lectures)}} \\ \hline
22 & Intermediate representations: DAGs, three-address code, quadruples/triples & CLO1, CLO3 & \S6.1--\S6.2 & [ET] \S5.1--5.3 \\ \hline
23 & Translating expressions and control-flow statements; backpatching & CLO2, CLO3 & \S6.4, \S6.6--\S6.7 & [AP] \S7.2--7.3 \\ \hline
24 & Type checking during translation; intermediate code for procedure calls & CLO1, CLO2 & \S6.3, \S6.5, \S6.9 & [AP] \S5.2, \S7.4 \\ \hline
25 & Building an intermediate-code generator: full AST $\to$ three-address-code translator \emph{(hands-on)} & CLO3 & \S6.1--\S6.9 (synthesis) & [ET] Ch.5 \\ \hline
\rowcolor{blue!8}
Opt.5 & \emph{\textbf{[Optional -- instructor's tool choice: LLVM, Part 1]}} Architecture, LLVM IR, key classes (Module, Function, BasicBlock, IRBuilder). Others may substitute GCC internals, a custom bytecode VM, or skip and use the time as buffer/review. & CLO3 & (tool choice for \S6.1--\S6.9) & [LLVM-LR] LLVM Language Reference Manual \\ \hline
\rowcolor{blue!8}
Opt.6 & \emph{\textbf{[Optional -- instructor's tool choice: LLVM, Part 2]}} Generating LLVM IR from an AST, end to end (hands-on). & CLO3 & (tool choice for \S6.1--\S6.9) & [LLVM-KA] LLVM Kaleidoscope Tutorial \\ \hline

\multicolumn{5}{|c|}{\cellcolor{gray!20}\textbf{Unit 6: Run-Time Environments (2 Lectures)}} \\ \hline
26 & Storage organization; stack allocation; activation records; nonlocal access & CLO1 & \S7.1--\S7.3 & [ET] Ch.6 \\ \hline
27 & Heap management; introduction to garbage collection & CLO1 & \S7.4--\S7.5 & [AP] Ch.13 \\ \hline

\multicolumn{5}{|c|}{\cellcolor{gray!20}\textbf{Unit 7: Code Generation \& Optimization (3 Lectures)}} \\ \hline
28 & Basic blocks \& flow graphs; local optimizations (constant folding, copy/constant propagation, CSE); peephole optimization & CLO1, CLO2 & \S8.4--\S8.5, \S8.7 & [MU] Ch.12 \\ \hline
29 & \textbf{Optimization Algorithm:} Common Subexpression Elimination via Available-Expressions data-flow analysis & CLO2 & \S9.1--\S9.3 & [MU] \S13.1--13.2 \\ \hline
30 & Register allocation \& assignment; instruction selection overview \emph{(previously-uncovered material from Ch.\,8, now restored)} & CLO1, CLO2 & \S8.8--\S8.9 & [AP] Ch.11 \\ \hline
\multicolumn{5}{|c|}{\cellcolor{yellow!25}\textit{--- Comprehensive Final Exam (covers L1--L30, emphasis on L21--L30) ---}} \\ \hline
\end{longtable}

\normalsize

\section*{Assessment Instruments}
\begin{itemize}
    \item \textbf{Project} -- students build a working compiler (lexer $\to$ parser $\to$ IR generator, optionally with basic optimization) for a small language, in stages aligned with the hands-on lectures (L7, L17, L25).
    \item \textbf{Quizzes} -- 7 short quizzes distributed across the semester (roughly one every 4 lectures), primarily on the lecture immediately preceding each quiz.
    \item \textbf{Assignments} -- 3--4 assignments: (1) regular expressions \& finite automata, (2) grammar transformation \& top-down parsing, (3) bottom-up parsing tables (SLR/CLR/LALR), and optionally (4) intermediate code generation / a chosen optimization.
    \item \textbf{Midterm 1} (after L10), \textbf{Midterm 2} (after L20), \textbf{Comprehensive Final} (after L30).
\end{itemize}

\section*{Tentative Grading Scheme}
\begin{itemize}
    \item Project -- 10 \%
    \item Assignments  -- 5 \%
    \item Quizzes -- 15 \%
    \item  Midterms -- 30 \%
    \item Final Exam -- 40 \%
\end{itemize}

\end{document}
